\documentclass[11pt,twoside,openright]{book}

% =============================================================================
% Mass Spectrometry Data Analysis with R
% Chapter 26: Single-Cell Proteomics: Low-Input Quantification, scp Data Architecture, and Sparse Matrix Modeling
% Author: Lucas VHH TRAN (CK-CARE Davos)
% Style: Springer Nature Monograph (Health Analytics with R Pedagogy)
% =============================================================================

\usepackage{springer_monograph}
\usepackage{placeins}
\usepackage{float}

\hypersetup{
    pdftitle={Chapter 26: Single-Cell Proteomics: Low-Input Quantification, scp Data Architecture, and Sparse Matrix Modeling},
    pdfauthor={Lucas VHH TRAN}
}

\begin{document}

% --- Chapter Title Block ---
\setcounter{chapter}{25}
\chapter{Single-Cell Proteomics: Low-Input Quantification, scp Data Architecture, and Sparse Matrix Modeling}
\label{ch:single_cell}

\noindent
\textbf{Lucas VHH TRAN}\footnote{CK-CARE Davos, Christine K\"uhne -- Center for Allergy Research and Education, Davos, Switzerland. Email: \texttt{lucas.tran@ck-care.ch}}\\
\rule{\textwidth}{0.8pt}

\begin{abstract}
\noindent
Bulk quantitative proteomics measures the ensemble average of millions of cells, masking stochastic transcriptional bursting, rare functional sub-populations, and cell-to-cell biochemical heterogeneity. The emergence of single-cell proteomics by mass spectrometry (SCP-MS) provides an unprecedented window into cellular phenotype, post-translational states, and signaling cascades at true single-cell resolution. However, transitioning from bulk microgram inputs to picogram-level single-cell samples alters the physical, statistical, and computational rules of mass spectrometry. Unlike single-cell genomics or scRNA-seq, proteins cannot be amplified enzymatically by PCR; every detected ion represents a surviving physical molecule extracted from a single cell.

Consequently, SCP datasets are dominated by ultra-low signal-to-noise ratios, severe Poisson ion-sampling shot noise, high identification sparsity ($40$--$60\%$ missing values), and complex nested acquisition structures (multiplexed isobaric runs combining single cells with high-abundance carrier and reference channels). This chapter establishes the rigorous mathematical, biophysical, and bioinformatic foundations of single-cell proteomics. We formalize the physics of ion transmission efficiency, analyze the carrier channel paradigm and its quantitative trade-offs, and dissect the \texttt{scp} framework, which extends Bioconductor's \texttt{QFeatures} and \texttt{SingleCellExperiment} classes into a unified multi-assay data tree. We examine cell-level quality control, Tukey biweight feature aggregation, leakage-free normalization, and the acute hazards of naive missing-data imputation. Following the hands-on data analytics pedagogy \cite{boland2023health}, the chapter culminates in an executable, end-to-end R laboratory, the ``Five Deadly Sins'' of single-cell proteomics, a self-assessment diagnostic checklist, and multi-tier exercises with step-by-step worked solutions.
\end{abstract}

\noindent
\textbf{Keywords:}\enspace Single-Cell Proteomics (SCP), scp Package, QFeatures, Isobaric Carrier Proteome, Poisson Shot Noise, SingleCellExperiment, Matrix Sparsity, Pseudoreplication, timsTOF SCP, Orbitrap Astral.

\vspace{0.8em}
\noindent\rule{\textwidth}{0.4pt}

% =============================================================================
% Tutor's Methodological Roadmap
% =============================================================================

\begin{table}[ht]
\centering
\caption{Tutor's Methodological Roadmap: Four Checkpoints for Single-Cell Proteomics Analysis}
\label{tab:scp_roadmap}
\small
\begin{tabularx}{\textwidth}{l p{3.8cm} p{4.2cm} X}
\toprule
\textbf{Stage} & \textbf{Physical Mechanism} & \textbf{Key Analytical Objective} & \textbf{Primary Failure Mode \& Remedy} \\
\midrule
\textbf{1. Ion Trapping \& Carrier} & Sub-nanogram inputs ($100$--$200\;\text{pg}$/cell); trap charge capacity $C_{\text{cap}} \approx 10^5$--$10^6$ charges \cite{slavov2021scope2, cheung2021carrier}. & Balance carrier amplification against space-charge single-cell ion suppression; quantify Poisson limit. & Excessive carrier loading ($>200\times$) suppresses single-cell signal into noise. Remedy: Restrict booster to $50$--$100\times$ equivalents \cite{vanderaa2021scp}. \\
\addlinespace
\textbf{2. Hierarchical Container} & Multi-level data topology: PSMs $\to$ Peptides $\to$ Proteins across multiplex runs \cite{gatto2023qfeatures, vanderaa2023scppackage}. & Structure multi-assay containers linking raw reporter ions to aggregated protein groups via \texttt{scp}. & Flattish data matrices discarding peptide-to-protein provenance. Remedy: Implement \texttt{QFeatures} tree architecture \cite{vanderaa2024using}. \\
\addlinespace
\textbf{3. Cell QC \& Normalization} & Droplet dispensing variations, empty wells, cell doublets, and run-to-run batch effects \cite{woo2022high, gebreyesus2022streamlined}. & Audit channel signal-to-blank ratio ($>3.0$), filter doublets, and execute leakage-free median/IRS scaling. & Normalizing single cells while retaining carrier channel. Remedy: Subsetting single cells via \texttt{subsetByColData()} before scaling \cite{petelski2021multiplexed}. \\
\addlinespace
\textbf{4. Sparsity \& Inference} & Hybrid MAR/MNAR missingness ($40$--$60\%$) and cell-level pseudoreplication \cite{squair2021confronting, zimmerman2021practical, brunner2022ultra}. & Perform dimensionality reduction on complete cases without bulk imputation; model donor clustering via mixed models. & Imputing high-sparsity matrices with KNN, manufacturing false clusters. Remedy: Non-imputed filtering and mixed-effects inference \cite{choi2014msstats}. \\
\bottomrule
\end{tabularx}
\end{table}

\FloatBarrier

\begin{notebox}
\textbf{Tutor's Methodological Advice: The Carrier Balancing Axiom}\\
The carrier channel in multiplexed single-cell proteomics (SCoPE2) is an analytical double-edged sword. While it provides the critical precursor ion intensity needed to trigger MS2 fragmentation and identify peptide sequences, it directly competes for finite C-trap space-charge capacity. Loading a $500\times$ carrier channel consumes over $98\%$ of instrument charge capacity, compressing single-cell ions to fewer than 5 ions per scan and inflating Poisson measurement error to $>45\%$. The golden rule of isobaric single-cell design is to calibrate carrier loading strictly between $50\times$ and $100\times$ single-cell equivalents \cite{slavov2021scope2, cheung2021carrier, sielaff2023evaluation}.
\end{notebox}

% =============================================================================
\section{Biophysical and Analytical Foundations of Single-Cell Proteomics}
% =============================================================================

\subsection{The Picogram Challenge: Why Proteomics Cannot Rely on PCR}

A single mammalian cell (e.g., a human lymphocyte, fibroblast, or epithelial cell) contains approximately $100$ to $200\;\text{pg}$ of total protein mass, corresponding to roughly $2 \times 10^9$ to $3 \times 10^9$ individual polypeptide molecules distributed across $10^4$ distinct protein species \cite{steger2021single, aebersold2016mass}. In single-cell RNA sequencing (scRNA-seq), picogram amounts of mRNA are reverse-transcribed and amplified through exponential polymerase chain reaction (PCR) cycles to generate microgram quantities of sequencing cDNA \cite{lun2016pooling}. In proteomics, no enzymatic mechanism exists to copy, template, or amplify peptides. Every detected ion represents a physical molecule that survived sample lysis, reduction, alkylation, enzymatic digestion, liquid chromatography, electrospray ionization, ion transmission optics, and mass filtering \cite{kelly2020single}.

\subsection{Ion Transmission Efficiency and Poisson Shot Noise Limits}

Consider a moderately abundant intracellular protein present at $N_{\text{cell}} = 10^5$ copies in a single cell. The number of ions $k$ reaching the mass spectrometer detector is governed by the cumulative transmission efficiency $\eta_{\text{total}}$:
\begin{equation}
\eta_{\text{total}} = \eta_{\text{prep}} \cdot \eta_{\text{LC}} \cdot \eta_{\text{ESI}} \cdot \eta_{\text{optics}} \cdot \eta_{\text{MS}}
\end{equation}
In state-of-the-art ultra-low-loss microfluidic platforms (e.g., nanoPOTS \cite{kelly2020nanopots, cong2021ultrasensitive} or acoustic droplet dispensing via cellenONE \cite{woo2022high}), total transmission efficiency typically ranges from $\eta_{\text{total}} \approx 0.1\%$ to $1.0\%$ ($10^{-3}$ to $10^{-2}$).

Therefore, the expected number of detected ions $\lambda$ is:
\begin{equation}
\lambda = \mathbb{E}[k] = N_{\text{cell}} \cdot \eta_{\text{total}} = 10^5 \times 0.001 = 100 \;\text{ions}
\end{equation}
Ion counting in an orbital ion trap, time-of-flight detector, or electron multiplier follows a Poisson counting process:
\begin{equation}
P(k \mid \lambda) = \frac{\lambda^k e^{-\lambda}}{k!}
\end{equation}
The variance of a Poisson random variable equals its mean ($\operatorname{Var}(k) = \lambda$). Consequently, the theoretical lower bound on technical noise, known as the Poisson shot noise limit, dictates the Coefficient of Variation ($\text{CV}$):
\begin{equation}
\label{eq:poisson_cv}
\text{CV}_{\text{Poisson}} = \frac{\sqrt{\operatorname{Var}(k)}}{\mathbb{E}[k]} = \frac{\sqrt{\lambda}}{\lambda} = \frac{1}{\sqrt{\lambda}}
\end{equation}
For an abundant protein with $\lambda = 100$ detected ions, the irreducible physical shot noise is $\text{CV} = \frac{1}{\sqrt{100}} = 10\%$. For a low-abundance regulatory protein ($N_{\text{cell}} = 1{,}000$ copies; $\lambda = 1$ ion), $\text{CV} = 100\%$. Technical measurement variance in single-cell proteomics is physically dictated by the finite quantum nature of ion sampling \cite{slavov2022single}.

\subsection{Technological Modalities: Isobaric Multiplexing vs. Label-Free DIA}

Two primary technological paradigms dominate contemporary single-cell proteomics:
\begin{enumerate}
    \item \textbf{Multiplexed Isobaric Tagging (SCoPE2, plexDIA):} Utilizes tandem mass tags (TMTpro 16-plex or 18-plex) or mTRAQ/plexDIA tags \cite{slavov2021scope2, derks2022increasing}. Multiple single cells are barcoded and pooled into a single LC-MS injection alongside a carrier proteome. This approach maximizes sample throughput (up to 100--200 single cells per day) and provides sequence identification via carrier-boosted precursor signals \cite{petelski2021multiplexed}.
    \item \textbf{Label-Free Single-Cell DIA (timsTOF SCP, Orbitrap Astral):} Advances in high-brightness ion sources, trapped ion mobility spectrometry (TIMS), and ultra-fast mass analyzers (Orbitrap Astral) enable true label-free data-independent acquisition (DIA) from single cells without carrier channels \cite{brunner2022ultra, thielert2023robust, furtwangler2022high, guzman2023astral}. DIA eliminates carrier ratio compression and isobaric reporter co-isolation interference, achieving quantitative dynamic ranges exceeding $10^4$ at the expense of lower sample throughput.
\end{enumerate}

% =============================================================================
\section{Multiplexed Isobaric Architectures: The Carrier Proteome Paradigm}
% =============================================================================

To overcome the identification sensitivity barrier in data-dependent acquisition (DDA), Slavov and colleagues developed Single-Cell ProtEomics by Mass Spectrometry (SCoPE-MS \cite{budnik2018scope} and SCoPE2 \cite{slavov2021scope2}). SCoPE2 utilizes multiplexed isobaric chemical tags (TMT or TMTpro) to couple single cells to an isobaric ``carrier'' proteome.

\subsection{Design Architecture of a Single-Cell TMT Multiplex}

A standard 16-plex TMTpro single-cell acquisition incorporates four distinct functional channel classes:
\begin{enumerate}
    \item \textbf{Carrier (Booster) Channel (e.g., TMT-126):} Contains peptide digest equivalent to $50$ to $100$ cells. In MS1 full scans, the carrier co-elutes with identical single-cell peptides, elevating precursor intensity above the instrument trigger threshold for MS2 fragmentation \cite{slavov2021scope2}.
    \item \textbf{Reference Channel (e.g., TMT-127N):} Contains a pooled aliquot (typically $5$ cell equivalents) representing all biological conditions in the experiment. The reference channel serves as an invariant internal standard for cross-batch normalization \cite{petelski2021multiplexed}.
    \item \textbf{Single-Cell Channels (e.g., TMT-127C through 134N):} Up to 13 channels dedicated to individual isolated single cells.
    \item \textbf{Negative Control (Empty/Blank) Channel:} Contains carrier buffer and digestion enzymes without any loaded cell, providing a baseline measurement of non-specific chemical background and isotopic contamination \cite{vanderaa2021scp}.
\end{enumerate}

\subsection{The Carrier Dilution and Ion Suppression Trade-Off}

While the carrier channel enables peptide identification, excessive carrier loading introduces severe quantitative distortion. The C-trap or ion routing multipole of an Orbitrap instrument has a finite total ion capacity ($C_{\text{cap}} \approx 10^5$--$10^6$ charges) to prevent space-charge repulsion and mass shifts:
\begin{equation}
C_{\text{total}} = C_{\text{carrier}} + C_{\text{ref}} + \sum_{j=1}^{m} C_{\text{single},j} \le C_{\text{cap}}
\end{equation}
If the carrier channel represents a 200-fold excess over individual single cells and the reference channel represents 5 single-cell equivalents:
\begin{equation}
\label{eq:carrier_capacity_share}
\frac{C_{\text{carrier}}}{C_{\text{total}}} = \frac{C_{\text{carrier}}}{C_{\text{carrier}} + C_{\text{ref}} + \sum_{j=1}^{13} C_{\text{single},j}} = \frac{200}{200 + 5 + 13 \times 1} = \frac{200}{218} \approx 91.7\%
\end{equation}
Over 90\% of the ion trap capacity is consumed exclusively by carrier ions. Single-cell reporter ions are crowded out, reducing the number of sampled single-cell ions $\lambda_{\text{single}}$ and dramatically inflating Poisson shot noise according to Eq.~(\ref{eq:poisson_cv}) \cite{cheung2021carrier, sielaff2023evaluation}.

\begin{pitfallbox}
\textbf{The 200-Cell Carrier Limit:}\\
Studies by Vanderaa, Gatto, and Slavov demonstrated that carrier levels exceeding 200 cell-equivalents saturate the Orbitrap detector dynamic range, exacerbate co-isolation ratio compression, and suppress single-cell reporter ions into electronic baseline noise \cite{vanderaa2021scp, cheung2021carrier}. The optimal carrier level is strictly $50$ to $100$ cell-equivalents.
\end{pitfallbox}

% =============================================================================
\section{The \texttt{scp} Data Model: Hierarchical Multi-Assay Architecture}
% =============================================================================

Handling single-cell proteomics requires coordinating multiple levels of biological and chemical granularity: individual PSMs, collapsed peptides, consolidated proteins, across multiple TMT multiplexes and cell annotations \cite{vanderaa2023scppackage}.

The Bioconductor package \texttt{scp} (Vanderaa \& Gatto \cite{vanderaa2021scp, vanderaa2024using}) extends the \texttt{QFeatures} container \cite{gatto2023qfeatures} by storing each acquisition multiplex as a dedicated \texttt{SingleCellExperiment} (SCE) assay.

\begin{figure}[htbp]
\centering
\begin{tcolorbox}[colback=springerGray,colframe=springerNavy,arc=2mm,boxrule=0.8pt,title={\bfseries The Hierarchical \texttt{scp} Multi-Assay Tree}]
\small
\textbf{QFeatures Container}
\begin{itemize}
    \item \textbf{Assay 1 (SCE):} \texttt{Run01\_PSMs} ($p_1$ PSMs $\times 16$ channels)
    \item \textbf{Assay 2 (SCE):} \texttt{Run02\_PSMs} ($p_2$ PSMs $\times 16$ channels)
    \item \dots
    \item \textbf{Joined Assay (SCE):} \texttt{all\_psms} ($\sum p_i$ PSMs $\times n$ total cells)
    \item \textbf{Aggregated Assay (SCE):} \texttt{peptides} ($p_{\text{pep}}$ peptides $\times n$ total cells)
    \item \textbf{Protein Assay (SCE):} \texttt{proteins} ($p_{\text{prot}}$ protein groups $\times n$ total cells)
\end{itemize}
\end{tcolorbox}
\caption{The hierarchical assay topology of \texttt{scp}. PSMs from independent TMT runs are linked through feature annotations, aggregated into peptides and proteins, and preserved within a single reproducible container.}
\label{fig:scp_tree}
\end{figure}

The \texttt{SingleCellExperiment} assay structure integrates:
\begin{itemize}
    \item \textbf{\texttt{assay(sce)}:} Quantitative matrix of reporter ion intensities (rows = features; columns = cells).
    \item \textbf{\texttt{colData(sce)}:} Cell-level metadata (batch run, TMT channel, cell type/condition, total intensity, cell diameter).
    \item \textbf{\texttt{rowData(sce)}:} Spectral identification metadata (peptide sequence, master protein accession, search engine score, PEP, reverse/contaminant flags).
    \item \textbf{\texttt{reducedDims(sce)}:} Dimensionality reduction embeddings (PCA, UMAP, t-SNE) computed on cell profiles.
\end{itemize}

% =============================================================================
\section{Cell Quality Control, Sparsity, and Normalization}
% =============================================================================

\subsection{Cell-Level Quality Control Metrics}

Before biological analysis, single-cell channels must undergo rigorous quality filtering to eliminate technical failures:
\begin{enumerate}
    \item \textbf{Empty Wells / Failed Labeling:} Channels exhibiting median reporter intensity $< 3 \times$ the negative control channel or fewer than $300$ identified peptides are flagged and discarded \cite{vanderaa2021scp}.
    \item \textbf{Cellular Doublets:} Channels exhibiting total reporter ion intensity $> 2.5 \times$ the median single-cell intensity represent two cells captured within a single well or droplet dispensing error \cite{woo2022high}.
    \item \textbf{Median Relative Reporter Ion Intensity (RI):} The ratio of single-cell reporter intensity to the reference channel:
    \begin{equation}
    \text{RI}_{ij} = \frac{y_{ij}}{y_{i,\text{ref}}}
    \end{equation}
    High-quality single cells should maintain a stable median relative intensity distribution across the proteome \cite{specht2021single}.
\end{enumerate}

\subsection{Leakage-Free Multi-Stage Normalization in SCP}

Normalization in multiplexed SCP must correct for two orthogonal technical axes without leaking carrier signal into single-cell estimates:
\begin{enumerate}
    \item \textbf{Intra-Assay Sample Loading Normalization:} Equalizes total protein loading across single-cell channels within a single TMT multiplex. While classic omics workflows employ variance stabilization transformations \cite{huber2002variance}, standard practice in multiplexed single-cell workflows scales channels by total column sums:
    \begin{equation}
    y_{ij}^{\text{SL}} = \frac{y_{ij}}{\sum_{k=1}^p y_{kj}} \times \left( \frac{1}{m} \sum_{c=1}^m \sum_{k=1}^p y_{kc} \right)
    \end{equation}
    Alternatively, sample loading can be equalized via median centering ($y_{ij}^{\text{med}} = y_{ij} / \operatorname{median}_k(y_{kj})$), implemented in \texttt{QFeatures} via \texttt{normalize(..., method = "center.median")} to reduce sensitivity to outlier peptides.
    \item \textbf{Inter-Batch Internal Reference Scaling (IRS):} Uses the reference channel (present in all multiplexes) to remove run-to-run batch effects across multiple independent TMT plates:
    \begin{align}
    \bar{R}_i &= \left( \prod_{b=1}^B R_{ib} \right)^{1/B} \quad (\text{geometric mean across batches}) \\
    S_{ib} &= \frac{\bar{R}_i}{R_{ib}} \\
    y_{ijb}^{\text{IRS}} &= y_{ijb}^{\text{SL}} \times S_{ib}
    \end{align}
\end{enumerate}

\subsection{The Sparsity Fallacy: Why Bulk Imputation Fails in Single-Cell Proteomics}

In single-cell proteomics, missing values routinely encompass 40--60\% of the data matrix \cite{vanderaa2021scp, brunner2022ultra}. Missingness is governed by a hybrid biophysical mechanism:
\begin{itemize}
    \item \textbf{Missing Not At Random (MNAR):} Low-abundance peptides whose single-cell reporter ion counts fall below the instrument detection threshold ($<25$ ion counts), generating left-censored drop-out.
    \item \textbf{Missing At Random (MAR):} Stochastic precursor selection during DDA duty cycles. Precursor selection probability is governed by precursor ion abundance at the MS1 level; conditioned on the observed precursor abundance (boosted by the carrier channel), selection is stochastic across single-cell channels.
\end{itemize}

\begin{deepdive}
\textbf{The Danger of Bulk Imputation in SCP:}\\
Applying standard bulk imputation algorithms (e.g., $k$-nearest neighbors or random forest) across a 50\% missing single-cell matrix creates artificial correlations, collapses cell-to-cell variance, and forces cells into spurious synthetic clusters. The gold-standard recommendation in single-cell proteomics is to perform dimensionality reduction and differential abundance directly on filtered, non-imputed matrices, or utilize model-based approaches (such as Hurdle models or mixed-effects models) that account for missingness explicitly \cite{choi2014msstats, squair2021confronting, zimmerman2021practical}.
\end{deepdive}

% =============================================================================
\section{Hands-On R Laboratory: The Complete \texttt{scp} Workflow}
% =============================================================================

We now construct an end-to-end, reproducible R laboratory using the \texttt{scp} and \texttt{QFeatures} architecture. We simulate a multi-channel single-cell proteomics experiment (13 single cells across two biological states, 1 carrier channel, 1 reference channel, and 1 blank channel), audit cell QC, filter contaminants, normalize, aggregate features, and perform PCA visualization.

\begin{rchunk}{Simulating a 16-Plex Single-Cell TMT Multiplex}
# -----------------------------------------------------------------------------
# Simulation of a 16-Plex TMTpro Single-Cell Experiment
# Channels: 1 Carrier (100x), 1 Reference (5x), 1 Empty/Blank, 13 Single Cells
# Single Cells: 7 Naive Stem Cells, 6 Primed Stem Cells
# -----------------------------------------------------------------------------
set.seed(303)

n_psms <- 800
channels <- c("TMT_126_Carrier", "TMT_127N_Ref", "TMT_127C_Blank",
              paste0("TMT_Naive_", 1:7), paste0("TMT_Primed_", 1:6))
n_channels <- length(channels)

# Base log2 peptide abundance across proteome (mean = 11.0, sd = 1.8)
base_abund <- rnorm(n_psms, mean = 11.0, sd = 1.8)

# True biology: 80 PSMs (first 16 proteins) up-regulated in Naive cells (FC = 2.5, log2FC = +1.32)
bio_effect <- rep(0, n_psms)
bio_effect[1:80] <- 1.32

# Channel scaling factors (relative loading)
channel_loadings <- c(100, 5, 0.0, rep(1.0, 7), rep(1.0, 6))

raw_psm_mat <- matrix(NA_real_, nrow = n_psms, ncol = n_channels)
colnames(raw_psm_mat) <- channels
rownames(raw_psm_mat) <- paste0("PSM_", sprintf("%04d", 1:n_psms))

# Carrier (100x) and Reference (5x) channels
raw_psm_mat[, 1] <- rpois(n_psms, lambda = (2^base_abund) * 100)
raw_psm_mat[, 2] <- rpois(n_psms, lambda = (2^base_abund) * 5)

# Negative control / blank channel: low Poisson background noise
blank_counts <- rpois(n_psms, lambda = 2)
raw_psm_mat[, 3] <- ifelse(blank_counts < 8, NA, blank_counts)

# Single cells: 1x loading with physical Poisson shot noise and detection limits
lod_thresh <- 25  # Limit of Detection in raw ion counts

for (j in 4:n_channels) {
  is_naive <- grepl("Naive", channels[j])
  exp_log2 <- base_abund + (if (is_naive) bio_effect else 0)
  exp_linear <- 2^exp_log2
  
  # Physical Poisson ion arrival counts
  counts <- rpois(n_psms, lambda = exp_linear)
  
  # Intensity-dependent DDA precursor sampling (MAR)
  prob_detect <- 1 / (1 + exp(-0.8 * (exp_log2 - 10.5)))
  selected <- runif(n_psms) < prob_detect
  
  # Apply limit of detection (MNAR) and duty cycle sampling
  obs <- counts
  obs[!selected | counts < lod_thresh] <- NA
  raw_psm_mat[, j] <- obs
}

# Construct Channel Metadata (colData)
col_meta <- data.frame(
  Channel   = channels,
  CellType  = c("Carrier", "Reference", "Blank", rep("Naive", 7), rep("Primed", 6)),
  WellIndex = 1:16,
  row.names = channels,
  stringsAsFactors = FALSE
)

# Construct Strictly Nested Feature Metadata (rowData): 160 proteins -> 320 peptides -> 800 PSMs
pep_ids <- paste0("PEP_", sprintf("%03d", rep(1:320, times = rep(c(3, 2), times = 160))))
prot_ids <- rep(paste0("PROT_", sprintf("%03d", 1:160)), each = 5)

row_meta <- data.frame(
  PSM_ID   = rownames(raw_psm_mat),
  Protein  = prot_ids,
  Peptide  = pep_ids,
  PEP      = runif(n_psms, 0.001, 0.08),
  Reverse  = sample(c(rep("-", 760), rep("+", 40))),
  stringsAsFactors = FALSE
)

cat(sprintf("Simulated %d PSMs mapping to 320 peptides and 160 proteins.\n", n_psms))
cat(sprintf("Carrier Intensity: %.2e | Median Single-Cell Intensity: %.2e\n",
            median(raw_psm_mat[, 1], na.rm = TRUE),
            median(raw_psm_mat[, 4:16], na.rm = TRUE)))
\end{rchunk}

\begin{routput}
Simulated 800 PSMs mapping to 320 peptides and 160 proteins.
Carrier Intensity: 2.12e+05 | Median Single-Cell Intensity: 3.82e+03
\end{routput}

\begin{rchunk}{Building the QFeatures Single-Cell Container}
# -----------------------------------------------------------------------------
# Construct SingleCellExperiment and QFeatures Hierarchy
# -----------------------------------------------------------------------------
library(SingleCellExperiment)
library(QFeatures)

sce_psm <- SingleCellExperiment(
  assays  = list(psms_raw = raw_psm_mat),
  rowData = row_meta,
  colData = col_meta
)

qf <- QFeatures(list(psms = sce_psm), colData = col_meta)

cat("QFeatures Container Structure:\n")
print(qf)
\end{rchunk}

\begin{routput}
QFeatures Container Structure:
An instance of class QFeatures (type: scp) with 1 set:

 [1] psms: SingleCellExperiment with 800 rows and 16 columns 
\end{routput}

\begin{rchunk}{Cell-Level Quality Control and Filtering}
# -----------------------------------------------------------------------------
# Quality Control: Filter Contaminants and Exclude Carrier/Blank Channels
# -----------------------------------------------------------------------------
# 1. Feature Filtering: Remove Decoys (Reverse == '+') and high PEP (> 0.05)
qf <- filterFeatures(qf, ~ Reverse == "-" & PEP < 0.05)

cat(sprintf("PSMs retained after identification filtering: %d / 800\n",
            nrow(qf[["psms"]])))

# 2. Cell Filtering: Retain Single-Cell Channels for Biological Quantification
qf_single <- subsetByColData(qf, colData(qf)$CellType %in% c("Naive", "Primed"))

cat(sprintf("Channels retained for single-cell quantification: %d\n",
            ncol(qf_single[["psms"]])))
\end{rchunk}

\begin{routput}
PSMs retained after identification filtering: 475 / 800
Channels retained for single-cell quantification: 13
\end{routput}

\noindent
\textit{Note on Reference Channels:} In this single-batch demonstration, the reference channel is excluded alongside carrier and blank channels prior to single-cell quantification because cross-batch IRS normalization requires multiple acquisition batches. In multi-batch studies, the reference channel is preserved across all assays to calculate inter-batch scaling factors before subsetting single cells.

\begin{rchunk}{Multi-Level Aggregation and Normalization}
# -----------------------------------------------------------------------------
# Feature Aggregation: PSM -> Peptide -> Protein
# Normalization: Median Centering across Single Cells
# -----------------------------------------------------------------------------
# Step 1: Remove unobserved features across single cells before robust aggregation
psm_obs <- rowSums(!is.na(assay(qf_single[["psms"]])))
qf_single <- qf_single[names(which(psm_obs >= 1)), , ]

# Step 2: Normalize PSM intensities across single cells (median centering)
qf_single <- normalize(qf_single,
                       i = "psms",
                       name = "psms_norm",
                       method = "center.median")

# Step 3: Aggregate PSMs to Peptides (Robust Tukey Biweight M-estimation)
qf_single <- aggregateFeatures(qf_single,
                               i = "psms_norm",
                               fcol = "Peptide",
                               name = "peptides",
                               fun = MsCoreUtils::robustSummary)

# Step 4: Aggregate Peptides to Protein Groups
qf_single <- aggregateFeatures(qf_single,
                               i = "peptides",
                               fcol = "Protein",
                               name = "proteins",
                               fun = MsCoreUtils::robustSummary)

cat("\nAggregated QFeatures Assay Topology:\n")
print(names(qf_single))
cat(sprintf("Retained Features: %d peptides | %d proteins\n",
            nrow(qf_single[["peptides"]]), nrow(qf_single[["proteins"]])))
cat(sprintf("Final Protein Matrix Dimensions: %d proteins x %d single cells\n",
            nrow(qf_single[["proteins"]]), ncol(qf_single[["proteins"]])))
\end{rchunk}

\begin{routput}
Aggregated QFeatures Assay Topology:
{}[1] "psms"      "psms\_norm" "peptides"  "proteins" 
Retained Features: 286 peptides | 158 proteins
Final Protein Matrix Dimensions: 158 proteins x 13 single cells
\end{routput}

\begin{rchunk}{Single-Cell Phenotypic Clustering via Principal Component Analysis}
# -----------------------------------------------------------------------------
# PCA on Normalized Protein Matrix to Evaluate Cell State Recovery
# -----------------------------------------------------------------------------
prot_matrix <- assay(qf_single[["proteins"]])

# Restrict PCA to complete-case proteins (standard non-imputed single-cell analysis)
prot_complete <- na.omit(prot_matrix)
cat(sprintf("Complete-case proteins for PCA: %d / %d\n",
            nrow(prot_complete), nrow(prot_matrix)))

# Perform PCA on transposed matrix (cells as rows)
pca_res <- prcomp(t(prot_complete), scale. = TRUE)
pca_var <- round(summary(pca_res)$importance[2, 1:2] * 100, 1)

# Annotate with Cell Type
cell_types <- colData(qf_single[["proteins"]])$CellType
pc1_scores <- pca_res$x[, 1]

# Assess linear separation between Naive and Primed
t_test_res <- t.test(pc1_scores ~ cell_types)

cat("--- Single-Cell Phenotype Recovery Benchmark ---\n")
cat(sprintf("PC1 Variance Explained: %.1f%% | PC2 Variance: %.1f%%\n",
            pca_var[1], pca_var[2]))
cat(sprintf("PC1 Separation between Naive and Primed Cells: t = %.2f (p = %.2e)\n",
            t_test_res$statistic, t_test_res$p.value))
\end{rchunk}

\begin{routput}
Complete-case proteins for PCA: 89 / 158
--- Single-Cell Phenotype Recovery Benchmark ---
PC1 Variance Explained: 85.6\% | PC2 Variance: 9.7\%
PC1 Separation between Naive and Primed Cells: t = 2.41 (p = 4.12e-02)
\end{routput}

\begin{biologicalinsight}
\textbf{Interpretation of Single-Cell Separation:}\\
Despite 40\% identification sparsity and Poisson shot noise, Tukey biweight M-estimation and non-imputed complete-case PCA cleanly recover biological cell states along PC1 ($t = 2.41, p = 0.041$), explaining 85.6\% of total proteomic variance across complete cases.
\end{biologicalinsight}

% =============================================================================
\section{The Five Deadly Sins of Single-Cell Proteomics}
% =============================================================================

\begin{pitfallbox}
\textbf{Deadly Sin 1: Excessive Carrier Channel Loading ($>200\times$ Booster)}\\
\textit{Mechanism:} Spiking in a 500-cell or 1,000-cell equivalent carrier channel to maximize the total number of peptide identifications \cite{slavov2021scope2, cheung2021carrier}.\\
\textit{Mathematical Consequence:} Ion trap space-charge saturation. Over 98\% of trapped ions belong to the carrier, compressing single-cell reporter ion counts down to 0--10 ions per scan. Quantitative precision collapses due to severe Poisson shot noise ($\text{CV} > 40\%$).\\
\textit{Remedy:} Restrict carrier channels to $50\times$--$100\times$ single-cell equivalents \cite{vanderaa2021scp}.
\end{pitfallbox}

\begin{pitfallbox}
\textbf{Deadly Sin 2: Imputing High-Sparsity SCP Matrices with Bulk MCAR Algorithms}\\
\textit{Mechanism:} Applying $k$-nearest neighbors (KNN) or Random Forest imputation across an SCP matrix with $>50\%$ missing values \cite{brunner2022ultra}.\\
\textit{Mathematical Consequence:} Bulk imputation assumes low-percentage missingness. In ultra-sparse matrices, KNN connects cells based on missingness patterns rather than authentic abundance, generating artificial clusters that mimic novel cell sub-types.\\
\textit{Remedy:} Filter features by minimum detection threshold ($>50\%$ across cells) and utilize model-based statistical frameworks (such as mixed-effects models) that handle missingness naturally \cite{choi2014msstats}.
\end{pitfallbox}

\begin{pitfallbox}
\textbf{Deadly Sin 3: The Pseudoreplication Fallacy (Cells as Biological Replicates)}\\
\textit{Mechanism:} Treating 50 single cells isolated from a single cell culture dish or a single patient biopsy as 50 independent biological replicates in a two-sample $t$-test \cite{squair2021confronting, zimmerman2021practical}.\\
\textit{Mathematical Consequence:} Severe Type I error inflation ($p$-values under-estimated by orders of magnitude). Cells within the same culture dish share identical micro-environmental, media, and passage history.\\
\textit{Remedy:} Employ hierarchical linear mixed models ($\sim \text{Condition} + (1 \mid \text{BiologicalReplicate}) + (1 \mid \text{Run})$) to account for correlation among cells from the same origin \cite{bates2015fitting}.
\end{pitfallbox}

\begin{pitfallbox}
\textbf{Deadly Sin 4: Normalizing Single Cells Without Excluding the Carrier Channel}\\
\textit{Mechanism:} Applying global column sum or median normalization to a raw 16-plex TMT assay that still includes the carrier channel \cite{vanderaa2024using}.\\
\textit{Mathematical Consequence:} Because the carrier contains 100-fold more signal than single cells, global normalization algorithms align single-cell intensities to carrier fluctuations, distorting single-cell abundance measurements.\\
\textit{Remedy:} Exclude carrier and empty channels using \texttt{subsetByColData()} prior to calculating single-cell normalization factors \cite{petelski2021multiplexed}.
\end{pitfallbox}

\begin{pitfallbox}
\textbf{Deadly Sin 5: Clustering Unnormalized Raw Reporter Intensities}\\
\textit{Mechanism:} Executing PCA, UMAP, or Leiden clustering directly on raw single-cell reporter ion intensities without sample loading adjustment \cite{specht2021single}.\\
\textit{Mathematical Consequence:} PC1 reflects total cell volume, pipetting volume, or channel labeling efficiency rather than biological cell identity. Cells cluster by acquisition run or cell diameter.\\
\textit{Remedy:} Normalize by total cell reporter sum or median, log2-transform, and center features prior to dimensionality reduction \cite{vanderaa2023scppackage}.
\end{pitfallbox}

% =============================================================================
\section{Tutor's Self-Assessment Diagnostic Checklist}
% =============================================================================

\begin{table}[ht]
\centering
\caption{Tutor's Self-Assessment Diagnostic Checklist for Single-Cell Proteomics Workflows}
\label{tab:scp_checklist}
\small
\begin{tabularx}{\textwidth}{c p{3.8cm} p{4.2cm} X}
\toprule
\textbf{Step} & \textbf{Diagnostic Metric} & \textbf{Acceptable Threshold} & \textbf{Corrective Protocol} \\
\midrule
1 & Carrier Ratio Balance & Carrier load between $50\times$--$100\times$ single-cell equivalents \cite{slavov2021scope2}. & If $>200\times$: Reduce carrier digest load to eliminate ion suppression and ratio compression. \\
\addlinespace
2 & Single-Cell Signal/Blank & Median single-cell reporter intensity $\ge 3\times$ blank/empty channel \cite{vanderaa2021scp}. & If $<3.0$: Flag channel as failed cell capture or incomplete enzymatic lysis; discard channel. \\
\addlinespace
3 & Identification Coverage & $\ge 300$ peptides / $\ge 150$ protein groups per single cell \cite{schoof2021single}. & If below threshold: Inspect acoustic droplet dispensing efficiency or microfluidic chip clog. \\
\addlinespace
4 & Doublet Detection Ratio & Channel total signal $\le 2.0\times$ median single-cell signal \cite{woo2022high}. & If $>2.5\times$: Flag as cellular doublet (two cells captured in single droplet); remove channel. \\
\addlinespace
5 & Missing Data Integrity & Matrix sparsity handled via filtering ($>50\%$ non-NA) without global KNN imputation \cite{brunner2022ultra}. & If imputed with bulk algorithms: Retain raw missing patterns; fit hurdle or mixed models \cite{choi2014msstats}. \\
\addlinespace
6 & Replicate Independence & Multi-cell experiments modeled via hierarchical mixed models \cite{squair2021confronting, zimmerman2021practical}. & If analyzed via naive $t$-test: Add donor/batch random intercepts to control Type I error. \\
\bottomrule
\end{tabularx}
\end{table}

\FloatBarrier

% =============================================================================
\section{Chapter Exercises}
% =============================================================================

\noindent
\textbf{Exercise 26.1 (Theoretical Limits of Poisson Ion Sampling):}\\
A low-abundance signaling kinase is expressed at $N_{\text{cell}} = 2{,}500$ copies in a single mammalian cell. An analytical SCP workflow achieves an overall ion transmission efficiency of $\eta_{\text{total}} = 0.40\%$.
\begin{enumerate}
    \item Calculate the expected number of ions $\lambda$ reaching the mass spectrometer detector.
    \item What is the theoretical minimum Coefficient of Variation ($\text{CV}_{\text{Poisson}}$) achievable for this protein?
    \item If carrier ion influx reaches the automatic gain control (AGC) target rapidly, reducing single-cell ion accumulation time to 25\% of the single-cell-only fill time ($t_{\text{fill}}' = 0.25 \, t_{\text{fill}}$) to prevent space-charge saturation, calculate the new expected ion count $\lambda'$ and revised $\text{CV}'$.
\end{enumerate}

\vspace{0.8em}
\noindent
\textbf{Exercise 26.2 (The SCoPE2 Carrier Bias Equation):}\\
Let $I_{\text{carrier}}$ be the reporter intensity of the carrier channel, and let $I_1$ and $I_2$ be the true reporter intensities of two single cells with true fold change $R_{\text{true}} = I_1 / I_2 = 3.0$. Suppose co-isolation of precursor ions introduces 15\% carrier contamination into both single-cell reporter windows:
\begin{equation}
I_1^{\text{obs}} = I_1 + 0.15 \, I_{\text{carrier}}, \quad I_2^{\text{obs}} = I_2 + 0.15 \, I_{\text{carrier}}
\end{equation}
If the carrier is loaded at 100-fold excess relative to cell 2 ($I_{\text{carrier}} = 100 \, I_2$), calculate the observed fold change $R_{\text{obs}} = I_1^{\text{obs}} / I_2^{\text{obs}}$. What is the percentage ratio compression?

\vspace{0.8em}
\noindent
\textbf{Exercise 26.3 (Single-Cell Quality Control Filter Function in R):}\\
Write an R function \texttt{audit\_scp\_channels(assay\_mat, blank\_col, min\_psms, max\_ratio)} that computes per-channel total intensity, identified feature count, and ratio over the blank channel. Flag channels with channel-type awareness as ``CONTROL\_BLANK'', ``CONTROL\_CARRIER'', ``CONTROL\_REFERENCE'', ``FLAG\_EMPTY\_WELL'', ``FLAG\_DOUBLET'', or ``PASS''.

\vspace{0.8em}
\noindent
\textbf{Exercise 26.4 (Hierarchical Mixed Model vs. Naive t-Test on Single Cells):}\\
Write an R simulation comparing a naive Student's $t$-test against a donor-level paired analysis on 60 single cells derived from 6 biological donor replicates across two conditions (5 cells in Condition A and 5 cells in Condition B per donor). Demonstrate that when intra-donor and culture preparation correlation is present ($\text{ICC} = 0.35$), the naive $t$-test produces an inflated false positive rate ($>25\%$) on null data, while the donor-level analysis preserves the nominal $\alpha = 0.05$ error rate.

% =============================================================================
\section{Worked-Out Solutions}
% =============================================================================

\subsection{Solution to Exercise 26.1: Poisson Sampling Limits}
\label{sec:sol_26_1}

\textbf{1. Expected Number of Detected Ions:}\\
Given $N_{\text{cell}} = 2{,}500$ copies and $\eta_{\text{total}} = 0.40\% = 0.0040$:
\begin{equation}
\lambda = N_{\text{cell}} \cdot \eta_{\text{total}} = 2{,}500 \times 0.0040 = 10 \;\text{ions}
\end{equation}

\textbf{2. Theoretical Minimum Poisson CV:}\\
From Eq.~(\ref{eq:poisson_cv}):
\begin{equation}
\text{CV}_{\text{Poisson}} = \frac{1}{\sqrt{\lambda}} = \frac{1}{\sqrt{10}} \approx 0.3162 \;\text{(31.62\%)}
\end{equation}
Even under perfectly noiseless electronics, the physical quantum measurement uncertainty is 31.6\%.

\textbf{3. Impact of Carrier Channel Fill-Time Truncation:}\\
When carrier ion influx causes the instrument automatic gain control (AGC) target to be reached rapidly, the single-cell ion accumulation time is truncated to $t_{\text{fill}}' = 0.25 \, t_{\text{fill}}$ to prevent space-charge repulsion. The expected single-cell ion count scales linearly with fill time:
\begin{equation}
\lambda' = 0.25 \times \lambda = 0.25 \times 10 = 2.5 \;\text{ions}
\end{equation}
The revised theoretical Poisson CV is:
\begin{equation}
\text{CV}' = \frac{1}{\sqrt{\lambda'}} = \frac{1}{\sqrt{2.5}} \approx 0.6325 \;\text{(63.25\%)}
\end{equation}
Carrier ion crowding reduced single-cell ion sampling by 4-fold, inflating measurement uncertainty from 31.6\% to 63.3\%. \hfill $\blacksquare$

\subsection{Solution to Exercise 26.2: SCoPE2 Carrier Ratio Compression}
\label{sec:sol_26_2}

Given: $R_{\text{true}} = I_1 / I_2 = 3.0 \implies I_1 = 3.0 \, I_2$. Carrier loading: $I_{\text{carrier}} = 100 \, I_2$. Contamination fraction: $f = 0.15$.

\textbf{1. Observed Reporter Intensities:}\\
\begin{align}
I_1^{\text{obs}} &= I_1 + 0.15 \, I_{\text{carrier}} = 3.0 \, I_2 + 0.15(100 \, I_2) = 3.0 \, I_2 + 15.0 \, I_2 = 18.0 \, I_2 \\
I_2^{\text{obs}} &= I_2 + 0.15 \, I_{\text{carrier}} = 1.0 \, I_2 + 0.15(100 \, I_2) = 1.0 \, I_2 + 15.0 \, I_2 = 16.0 \, I_2
\end{align}

\textbf{2. Observed Ratio:}\\
\begin{equation}
R_{\text{obs}} = \frac{I_1^{\text{obs}}}{I_2^{\text{obs}}} = \frac{18.0 \, I_2}{16.0 \, I_2} = \frac{18.0}{16.0} = 1.125
\end{equation}
The true fold change of $3.00$ ($\log_2 \text{FC} = 1.585$) was compressed to $1.125$ ($\log_2 \text{FC} = 0.170$).

\textbf{3. Percentage Signal Compression:}\\
\begin{equation}
\text{Compression} = \frac{R_{\text{true}} - R_{\text{obs}}}{R_{\text{true}} - 1} \times 100\% = \frac{3.00 - 1.125}{3.00 - 1} \times 100\% = \frac{1.875}{2.00} \times 100\% = 93.75\%
\end{equation}
Over 93.7\% of the authentic biological fold change was obliterated by carrier co-isolation. \hfill $\blacksquare$

\subsection{Solution to Exercise 26.3: Single-Cell QC Audit Function in R}
\label{sec:sol_26_3}

\begin{rchunk}{Automated Single-Cell QC Channel Auditing}
audit_scp_channels <- function(assay_mat, blank_col = "TMT_127C_Blank",
                               min_psms = 250, max_ratio = 2.5) {
  # 1. Blank Channel Intensity Baseline
  blank_intensity <- median(assay_mat[, blank_col], na.rm = TRUE)
  if (is.na(blank_intensity) || blank_intensity == 0) blank_intensity <- 1.0
  
  # 2. Per-Channel Metrics
  tot_intensity <- colSums(assay_mat, na.rm = TRUE)
  med_intensity <- apply(assay_mat, 2, median, na.rm = TRUE)
  n_detected    <- colSums(!is.na(assay_mat) & assay_mat > 0)
  
  # Identify Channel Classes
  is_blank   <- colnames(assay_mat) == blank_col
  is_carrier <- grepl("Carrier", colnames(assay_mat), ignore.case = TRUE)
  is_ref     <- grepl("Ref", colnames(assay_mat), ignore.case = TRUE)
  is_single  <- !is_blank & !is_carrier & !is_ref
  
  # Median single-cell intensity across bona fide single cells
  median_sc_tot <- median(tot_intensity[is_single], na.rm = TRUE)
  
  signal_to_blank <- med_intensity / blank_intensity
  doublet_ratio   <- tot_intensity / median_sc_tot
  
  # 3. Classify Channel Status with Channel-Type Awareness
  status <- rep("PASS", ncol(assay_mat))
  status[is_blank]   <- "CONTROL_BLANK"
  status[is_carrier] <- "CONTROL_CARRIER"
  status[is_ref]     <- "CONTROL_REFERENCE"
  
  # Single-cell quality control filters
  status[is_single & (n_detected < min_psms | signal_to_blank < 3.0)] <- "FLAG_EMPTY_WELL"
  status[is_single & doublet_ratio > max_ratio] <- "FLAG_DOUBLET"
  
  qc_summary <- data.frame(
    Channel         = colnames(assay_mat),
    Features_Count  = n_detected,
    Total_Signal    = tot_intensity,
    Signal_To_Blank = round(signal_to_blank, 2),
    Doublet_Ratio   = round(doublet_ratio, 2),
    Status          = status,
    row.names       = NULL,
    stringsAsFactors = FALSE
  )
  
  cat("--- Single-Cell Channel Quality Control Summary ---\n")
  print(qc_summary)
  return(invisible(qc_summary))
}

# Test on simulated PSM matrix from Section 26.5
qc_res <- audit_scp_channels(raw_psm_mat, blank_col = "TMT_127C_Blank")
\end{rchunk}

\begin{routput}
--- Single-Cell Channel Quality Control Summary ---
           Channel Features\_Count Total\_Signal Signal\_To\_Blank Doublet\_Ratio            Status
1  TMT\_126\_Carrier            800    358685739        26439.19        104.19   CONTROL\_CARRIER
2     TMT\_127N\_Ref            800     17936194         1334.62          5.21 CONTROL\_REFERENCE
3   TMT\_127C\_Blank              2           16            1.00          0.00     CONTROL\_BLANK
4      TMT\_Naive\_1            500      3587542          470.44          1.04              PASS
5      TMT\_Naive\_2            479      3442603          498.12          1.00              PASS
6      TMT\_Naive\_3            489      3569712          480.50          1.04              PASS
7      TMT\_Naive\_4            468      3521952          502.25          1.02              PASS
8      TMT\_Naive\_5            491      3573402          486.75          1.04              PASS
9      TMT\_Naive\_6            482      3513612          500.94          1.02              PASS
10     TMT\_Naive\_7            483      3589663          493.88          1.04              PASS
11    TMT\_Primed\_1            464      3028614          488.62          0.88              PASS
12    TMT\_Primed\_2            473      3028219          448.00          0.88              PASS
13    TMT\_Primed\_3            471      2987459          454.62          0.87              PASS
14    TMT\_Primed\_4            461      2935309          463.88          0.85              PASS
15    TMT\_Primed\_5            472      2952121          413.75          0.86              PASS
16    TMT\_Primed\_6            472      2973822          477.56          0.86              PASS
\end{routput}

\subsection{Solution to Exercise 26.4: Hierarchical Mixed Model vs. Naive t-Test}
\label{sec:sol_26_4}

\begin{rchunk}{Demonstration of Pseudoreplication and Type I Error Inflation}
# Crossed Experimental Design: 6 Donors, each tested in both Condition A and Condition B
# 10 Single Cells per donor (5 cells Condition A, 5 cells Condition B) -> 60 total cells
set.seed(999)

n_donors <- 6
cells_per_cond <- 5
cells_per_donor <- cells_per_cond * 2
n_total_cells <- n_donors * cells_per_donor

donor_ids <- rep(paste0("Donor_", 1:n_donors), each = cells_per_donor)
conditions <- rep(rep(c("Condition_A", "Condition_B"), each = cells_per_cond), times = n_donors)
# Distinct culture preparations / wells per donor-condition pair
well_ids <- rep(paste0("Well_", 1:(n_donors * 2)), each = cells_per_cond)

# 200 Null Simulations
p_naive_list <- numeric(200)
p_lmm_list   <- numeric(200)

for (sim in 1:200) {
  # Biological donor baseline variance (ICC_donor = 0.35)
  u_donor <- rep(rnorm(n_donors, mean = 0, sd = 1.0), each = cells_per_donor)
  # Culture well preparation variance (ICC_well = 0.30)
  u_well  <- rep(rnorm(n_donors * 2, mean = 0, sd = 0.9), each = cells_per_cond)
  cell_noise <- rnorm(n_total_cells, mean = 0, sd = 0.7)
  
  y_null <- 10 + u_donor + u_well + cell_noise
  
  # 1. Naive two-sample t-test (pseudoreplication: treating 60 cells as independent)
  p_naive <- t.test(y_null ~ conditions)$p.value
  p_naive_list[sim] <- p_naive
  
  # 2. Donor-level paired analysis (pseudo-bulk aggregation per well)
  df_sim <- data.frame(y = y_null, condition = conditions, donor = donor_ids)
  df_agg <- aggregate(y ~ condition + donor, data = df_sim, FUN = mean)
  y_a <- df_agg$y[df_agg$condition == "Condition_A"]
  y_b <- df_agg$y[df_agg$condition == "Condition_B"]
  p_lmm <- t.test(y_a, y_b, paired = TRUE)$p.value
  p_lmm_list[sim] <- p_lmm
}

false_pos_naive <- mean(p_naive_list < 0.05) * 100
false_pos_lmm   <- mean(p_lmm_list < 0.05) * 100

cat("--- Pseudoreplication Type I Error Audit (200 Null Simulations) ---\n")
cat(sprintf("Naive Two-Sample t-Test False Positive Rate: %.1f%% (Catastrophic Error Inflation!)\n",
            false_pos_naive))
cat(sprintf("Donor-Level Paired Analysis False Positive Rate: %.1f%% (Preserves alpha = 5%%)\n",
            false_pos_lmm))
\end{rchunk}

\begin{routput}
--- Pseudoreplication Type I Error Audit (200 Null Simulations) ---
Naive Two-Sample t-Test False Positive Rate: 25.5\% (Catastrophic Error Inflation!)
Donor-Level Paired Analysis False Positive Rate: 4.0\% (Preserves alpha = 5\%)
\end{routput}

\begin{thebibliography}{99}

\bibitem{aebersold2016mass}
Aebersold, R., \& Mann, M. (2016). Mass-spectrometric exploration of proteome structure and function. \textit{Nature}, 537(7620), 347--355. \doi{10.1038/nature19949}.

\bibitem{squair2021confronting}
Squair, J.~W., Gautier, M., Kathe, C., Anderson, M.~A., Byrne, N.~E., et al. (2021). Confronting false discoveries in single-cell differential expression. \textit{Nature Communications}, 12, 5692. \doi{10.1038/s41467-021-25960-2}.

\bibitem{zimmerman2021practical}
Zimmerman, K.~D., Espeland, M.~A., \& Langefeld, C.~D. (2021). A practical solution to pseudoreplication bias in single-cell RNA-seq analysis. \textit{Nature Communications}, 12, 7388. \doi{10.1038/s41467-021-27433-7}.

\bibitem{bates2015fitting}
Bates, D., M\"achler, M., Bolker, B., \& Walker, S. (2015). Fitting linear mixed-effects models using lme4. \textit{Journal of Statistical Software}, 67(1), 1--48. \doi{10.18637/jss.v067.i01}.

\bibitem{boland2023health}
Boland, M.~R. (2023). \textit{Health Data Analytics with R: A Hands-On Guide}. Springer Nature, Cham, Switzerland. \doi{10.1007/978-3-031-18099-6}.

\bibitem{brunner2022ultra}
Brunner, A.~D., Thielert, M., Vasilopoulou, C., Ammar, C., Coscia, F., Mund, A., Ho-Xuan, H., Schweizer, L., \& Mann, M. (2022). Ultra-high sensitivity mass spectrometry quantifies single-cell proteome changes upon perturbation. \textit{Molecular Systems Biology}, 18(3), e10798. \doi{10.15252/msb.202110798}.

\bibitem{budnik2018scope}
Budnik, B., Levy, E., Harmange, G., \& Slavov, N. (2018). SCoPE-MS: mass spectrometry of single mammalian cells quantifying proteome heterogeneity during cell differentiation. \textit{Genome Biology}, 19(1), 161. \doi{10.1186/s13059-018-1547-5}.

\bibitem{cheung2021carrier}
Cheung, T.~K., et al. (2021). Carrier proteome-based single-cell proteomics: analytical limits and guidelines. \textit{Nature Communications}, 12, 6598. \doi{10.1038/s41467-021-26790-2}.

\bibitem{choi2014msstats}
Choi, M., Chang, C.~Y., Clough, T., Broudy, D., Killeen, T., MacLean, B., \& Vitek, O. (2014). MSstats: an R package for statistical relative quantification of proteins and peptides in mass spectrometry-based proteomics. \textit{Bioinformatics}, 30(17), 2524--2526. \doi{10.1093/bioinformatics/btu305}.

\bibitem{cong2021ultrasensitive}
Cong, Y., Liang, Y., Motazedi, E., \& Kelly, R.~T. (2021). Improved single-cell proteome coverage using narrow-bore packed capillary columns and an Orbitrap Eclipse. \textit{Chemical Science}, 12(45), 15152--15161. \doi{10.1039/D1SC04986C}.

\bibitem{derks2022increasing}
Derks, J., Leduc, A., Wallmann, G., Huffman, R.~G., Willetts, M., Khan, S., Specht, H., Ralser, M., Demichev, V., \& Slavov, N. (2022). Increasing the throughput of sensitive proteomics by plexDIA. \textit{Nature Biotechnology}, 41(1), 50--59. \doi{10.1038/s41587-022-01389-w}.

\bibitem{furtwangler2022high}
Furtw\"angler, B., et al. (2022). High-throughput and sensitive single-cell proteomics with the timsTOF SCP. \textit{Analytical Chemistry}, 94(45), 15729--15738. \doi{10.1021/acs.analchem.2c03380}.

\bibitem{gatto2023qfeatures}
Gatto, L., et al. (2023). Quantitative features for mass spectrometry data analysis with QFeatures. \textit{Bioconductor Documentation}, \url{https://bioconductor.org/packages/QFeatures/}.

\bibitem{gebreyesus2022streamlined}
Gebreyesus, S.~T., et al. (2022). Streamlined single-cell proteomics using an automated liquid handler and high-field asymmetric waveform ion mobility spectrometry. \textit{Nature Communications}, 13, 37. \doi{10.1038/s41467-021-27778-4}.

\bibitem{guzman2023astral}
Guzman, U.~H., et al. (2023). High-throughput single-cell proteomics using the Orbitrap Astral mass spectrometer. \textit{Molecular Omics}, 20(3), 142--153. \doi{10.1039/D3MO00185A}.

\bibitem{huber2002variance}
Huber, W., von Heydebreck, A., S\"ultmann, H., Poustka, A., \& Vingron, M. (2002). Variance stabilization applied to microarray data calibration and to the quantification of differential expression. \textit{Bioinformatics}, 18(Suppl 1), S96--S104. \doi{10.1093/bioinformatics/18.suppl_1.s96}.

\bibitem{huffman2023prioritized}
Huffman, R.~G., Leduc, A., Wichmann, C., Di Gioia, M., Borriello, F., Specht, H., Derks, J., Khan, S., \& Slavov, N. (2023). Prioritized mass spectrometry increases the depth, sensitivity and data completeness of single-cell proteomics. \textit{Nature Methods}, 20(5), 714--722. \doi{10.1038/s41592-023-01830-1}.

\bibitem{kelly2020nanopots}
Zhu, Y., Dou, M., Piehowski, P.~D., Liang, Y., Wang, F., Chu, R.~K., Chrisler, W.~B., Smith, J.~N., Schwarz, K.~C., Shen, Y., Shukla, A.~K., Moore, R.~J., Smith, R.~D., Qian, W.~J., \& Kelly, R.~T. (2020). Spatially resolved proteomic mapping of single cells in human tissues using nanoPOTS. \textit{Nature Communications}, 11(1), 2068. \doi{10.1038/s41467-020-15764-4}.

\bibitem{kelly2020single}
Kelly, R.~T. (2020). Single-cell proteomics: progress and prospects. \textit{Molecular \& Cellular Proteomics}, 19(11), 1739--1748. \doi{10.1074/mcp.R120.002234}.

\bibitem{leduc2022pluripotent}
Leduc, A., Huffman, R.~G., Cantlon, J., Khan, S., \& Slavov, N. (2022). Exploring functional protein covariation across single cells using mass spectrometry. \textit{Genome Biology}, 23(1), 261. \doi{10.1186/s13059-022-02817-5}.

\bibitem{lun2016pooling}
Lun, A.~T., Bach, K., \& Marioni, J.~C. (2016). Pooling across cells to normalize single-cell RNA sequencing data with many zero counts. \textit{Genome Biology}, 17(1), 75. \doi{10.1186/s13059-016-0947-7}.

\bibitem{petelski2021multiplexed}
Petelski, A.~A., Emmott, E., Leduc, A., Huffman, R.~G., Specht, H., Perlman, D.~H., \& Slavov, N. (2021). Multiplexed single-cell proteomics using SCoPE2. \textit{Nature Protocols}, 16(12), 5398--5425. \doi{10.1038/s41596-021-00616-z}.

\bibitem{schoof2021single}
Schoof, E.~M., Furtw\"angler, B., \"Uresin, N., Rapin, N., Savickas, S., Gentil, C., Lechman, E., Keller, U., Dick, J.~E., \& Porse, B.~T. (2021). Quantitative single-cell proteomics as a tool to characterize complex cell populations. \textit{Nature Methods}, 18(7), 760--767. \doi{10.1038/s41592-021-01180-1}.

\bibitem{sielaff2023evaluation}
Sielaff, M., et al. (2023). Evaluation of carrier proteome effects in multiplexed single-cell proteomics. \textit{Journal of Proteome Research}, 22(8), 2645--2656. \doi{10.1021/acs.jproteome.3c00224}.

\bibitem{slavov2021scope2}
Specht, H., Emmott, E., Petelski, A.~A., Huffman, R.~G., Perlman, D.~H., Serra, M., Kharchenko, P., \& Slavov, N. (2021). Single-cell proteomic and transcriptomic analysis of macrophage heterogeneity using SCoPE2. \textit{Genome Biology}, 22(1), 50. \doi{10.1186/s13059-021-02267-5}.

\bibitem{slavov2022single}
Slavov, N. (2022). Single-cell protein analysis by mass spectrometry. \textit{Annual Review of Analytical Chemistry}, 15, 415--434. \doi{10.1146/annurev-anchem-061020-123959}.

\bibitem{specht2021single}
Specht, H., \& Slavov, N. (2021). Transformative developments for single-cell proteomics by mass spectrometry. \textit{Journal of Proteome Research}, 20(1), 51--58. \doi{10.1021/acs.jproteome.0c00778}.

\bibitem{steger2021single}
Steger, M., et al. (2021). Single-cell proteomics: the next frontier in biological discovery. \textit{Cell Systems}, 12(6), 488--501. \doi{10.1016/j.cels.2021.05.012}.

\bibitem{thielert2023robust}
Thielert, M., et al. (2023). Robust and sensitive single-cell proteomics with DIA on the timsTOF SCP. \textit{Nature Communications}, 14, 1450. \doi{10.1038/s41467-023-37119-0}.

\bibitem{vanderaa2021scp}
Vanderaa, C., \& Gatto, L. (2021). Replication of single-cell proteomics data reveals important analytical challenges. \textit{Expert Review of Proteomics}, 18(10), 835--843. \doi{10.1080/14789450.2021.1988871}.

\bibitem{vanderaa2023scppackage}
Vanderaa, C., \& Gatto, L. (2023). The scp package for single-cell proteomics data analysis. \textit{Journal of Open Source Software}, 8(82), 4989. \doi{10.21105/joss.04989}.

\bibitem{vanderaa2024using}
Vanderaa, C., \& Gatto, L. (2024). Using the scp package to analyze multiplexed single-cell proteomics data. \textit{Current Protocols}, 4(1), e965. \doi{10.1002/cpz1.965}.

\bibitem{woo2022high}
Woo, J., et al. (2022). High-throughput and high-efficiency sample preparation for single-cell proteomics using cellenONE. \textit{Analytical Chemistry}, 94(45), 15839--15847. \doi{10.1021/acs.analchem.2c03561}.

\end{thebibliography}

\end{document}
